\section*{A city choice built around a student's day}

Zurich is the cautious recommendation emerging from our comparison of sixteen cities for high-school life. An alternative that places greater emphasis on education points to Oslo. Together, these findings offer students and families a practical starting point: a city comparison that connects everyday opportunities with personal priorities while keeping gaps in the evidence visible.

For a teenager, choosing a city involves the school day and the hours around it. Teaching resources, nearby higher education, clean air, safety, housing and living costs, transport, journey times, internet access and swimming facilities can all matter. Our study brings these ten considerations into one framework. Its purpose is to help families examine trade-offs, rather than assume that the same place will suit every student.

\subsection*{A cautious recommendation with a clear meaning}

The comparison uses common scoring yardsticks and a stated balance of priorities. Where a figure is missing or describes part of a city's resources, the model keeps a range of possible scores. An unknown figure therefore remains open rather than being replaced with an invented average. Each city's overall result is a range, reflecting the information available under the model's assumptions.

The recommendation asks how far a selected city could fall behind a better-scoring alternative within those ranges. Zurich has the smallest worst-case shortfall under the baseline priorities. It also remains the cautious recommendation when all ten considerations receive equal weight. This makes it a useful place to begin a broad comparison; all sixteen cities nevertheless remain possible leaders within the retained evidence ranges. The study has not established a single certain winner.

\subsection*{Better evidence makes the comparison more useful}

Official records provide concrete improvements to the comparison. Zurich's transport operator reports 446 stops on city territory in its 2025 facts booklet; separate figures for the surrounding area are kept out of the city count. Stockholm's published school indicator measures pupils per full-time-equivalent teacher in upper-secondary schools located in the municipality, across providers. These examples show why geography, school stage and the meaning of a count matter when comparing cities.

Across fifteen previously missing entries, additional evidence reduced the average width of the city score ranges by about 14 per cent. The original submitted values and scoring rules stayed in place. Narrower ranges make the comparison more informative, while the full set of possible leaders shows that further evidence could still change the conclusion. The records span different years, and some describe school or city resources rather than the experience of a particular teenager.

\newpage
\subsection*{Students' priorities can change the recommendation}

In an education-focused comparison, the relative emphasis on teaching resources and the density of higher-education institutions is doubled before the priorities are balanced again. Oslo then becomes the cautious recommendation. This is an illustrative choice made by the analysts, not a finding from a student survey. It gives a family a concrete way to explore how a stronger educational emphasis changes its shortlist.

City totals also need a local check. A transport-stop inventory leaves questions about walking distance, service frequency, fares and accessible routes. A pool count leaves questions about entry fees, opening hours and which neighbourhoods the facilities serve. These are useful questions for families and city planners to investigate before treating resource supply as equal access. The study does not measure an individual student's benefit from adding a stop, teacher or pool.

\subsection*{Use future scenarios to ask better questions}

The study also examines three hypothetical changes: expanding selected resources, increasing living and travel pressures, and reducing selected educational, air-quality and travel burdens. When the affected quantities change uniformly by ten per cent, Zurich remains the cautious recommendation under both baseline and equal priorities. Repeated simulations explore what happens when the assumed changes differ across cities. They test chosen scenarios; they do not predict a future year or estimate a city's real chance of being best.

A family's next step is to pair the city comparison with its own school options, budget, daily route and preferred activities. Students can give education more emphasis, examine the affordability trade-offs, and check whether nearby opportunities fit their routines. The model makes those choices easier to discuss: a clear starting recommendation, room for different priorities, and an explicit view of the evidence needed for a more personal decision.
